% An interesting discussion I've been playing around with is a system of normative ethics via Golden reasoning. However, I should note that this system itself is not ready-for-use in practical circumstances, but rather is more of an intrigue class akin to worldbuilding or something in the sense that it would only be a genuinely useful practice in the case of having more of a god's-eye-vantage than we as humans actually have on demand. Specifically, the idea is this: ``Golden ethics becomes practically useful assuming you have on-hand the quantity of how each stated-attribute under a given Aethic attributes relatively affects the probability of the Golden Point under the Cosmic Nexus". Anything short of that though, and you're a bit out of luck, because the requirement for this mode of inference is actually having such available. But, the reason I think of this as like a kind of worldbuilding is because if you \textit{did} actually have this data on-hand on demand, you could design around it some interesting categorizations of things. So in that sense it's sort of like utilitarianism I guess---although hopefully better because it doesn't get tripped up by a few of its paradoxes---I don't think I wrote down the derivations for those so I have to re-derive them though lol once I read the standard utilitarianism paradoxes. But anyways I digress---the last thing I want to mention for right now is that the Golden ethics concept is indeed different from utilitarianism because it's all about conditioning on infinite-biosphere-survival rather than utiliy, and also I want to clearly state for the record that it isn't longtermism either, for a variety of reasons but one quick and easy one is that longtermism is just utilitarianism in a different dress, which also conditions on utility---which note is an optimization-type conditioning rather than the straight up more abstract Bayesian-style-but-for-boundary-conditions-style conditioning of Golden reasoning---(oh wait actually thinking about this and one of the arguments I was talking about just came back to me---Golden reasoning would rate 1000 people living farther into the future higher than a mllion times more people living to not as far in the future (but technically the rating's grading is probably ``probability of living to infinite boundary future" rather than ``how far into the future does anybody last" non-probabalistically), so it's about ``how far out can we get the last member of the biosphere to get" rather than naively adding population like longtermism! So emergently/indirectly it might help to boost population if that to probabalistic precision likely pushes the farthest person out further into the future, but even still if in such a scenario the population is decaying exponentially then that's not exactly arbitrarily-far sustainability, so that would rank poorly as Departric anyway! Anyways the point of the matter is that the paradox I'd heard was ``should we build a town on this region of land near industrial waste" would by longtermism be ranked ``yes", because all those people would die young, but making them exist is better than making them not exist---so that exact principle of longtermism is essentially nullified as a primtive in the Golden ethics sense, in the sense that Golden ethics ranges anywhere from uncorrelation---ie even advantage (so at least it doesn't positively say ``yea go build that town mfers")---to positive correlation in cases like the fate of the whole species). 

\section{Golden Ethics}

Golden Ethics is a hypothetical ethical system that is derived from the Golden Initiative. The basic premise is that if we had full knowledge of which actions make the Golden Point more or less likely, then Golden Ethics concerns what our most ideal actions as agents would be. Of course, we do not actually have access to such an omniscience about the Golden Point at present, however we can nonetheless gain theoretical ground by supposing how such a system of ethics would be constructed if the required information was present. That is, if we knew our goal, what would be our optimal course of action given that? Any system of information less ideal than this high point, then, we could imagine as a kind of approximation of the omniscient scenario to be discussed here. 

To begin with, let us state the fundamental principle of Golden Ethics. 
\begin{prince}[Fundamental Principle of Golden Ethics]
    One ought to consistently pick the decision with the highest Golden Scores when given the option. 
\end{prince}
\begin{definition}[Golden Virtue]
    The act of abiding by the fundamental principle of Golden Ethics. 
\end{definition}
\begin{definition}[The Golden Obligation]
    The obligation under Golden Ethics to act in a Goldenly Virtuous manner. 
\end{definition}
Note, also, what the Golden Ethical system of judgment ought to be, which is very different than the classical system. Suppose I am person $A$, and I want to choose whether or not I should judge person $B$ for an action they committed. Classically, we might suppose that the answer ought to be in terms of \textit{whether or not person $B$ acted ethically, so in this case abided by the fundamental principle}. But this, you see, is actually a mischaracterization of that principle. In Golden Ethics, the primary rule is to follow the fundamental principle, so this means that \textit{I should actually judge person $B$ based on whether or not my own act of judgment would be following the fundamental principle}. So the key difference from the classical ethical mindset is that judgment should always be based in whether it is ethical to judge, even more so than if the defendant acted ethically or not. It is trivial to see that these two notions should usually line up, however in those rare stances when they do not, the Golden Ethical version ought to be defaulted to. Let us lay out some intuitive examples of why this is so.

Consider a situation where a person commits an extremely antisocial offense by trying to ruin or even harm another person. However, in doing so, they accidentally provide their would-be-victim with a new avenue toward a breakthrough. Flukes like this happen all the time in high-Golden Score cases, simply due to the fundamental theorem of Nexusive reasoning, (such that the antagonist's actions actually work as a stepping stone for the protagonist across some boundary, with the classic example being theropod dinosaurs in the Dark Era -- in order to optimize our chances of gaining mammalian cognition, Optimation provided us with a scenario that was just hazardous enough to force evolution in us, but just safe enough that we had a shot, albeit minimal, of surviving through it. Note that this exact kind of stepping stone scenario is strikingly common in both natural history and frankly any system where the fundamental theorem of Nexusive reasoning applies). In such a scenario, they might very well provide you with many Golden Scores out of being helpful, however they have also done a horrendous deed. As such, we might intuitively suppose that it is more important for us to extrapolate their future deeds than it is for us to reward them for a single accidental good deed. This is essentially a consequentialist reformulation of the tenets of virtue ethics. Generally, we are relying on the low probability of a bad person repeating a good deed to form our judgment of them. This can be seen as a generalization of the tenets of virtue ethics, where we are updating its notion that a person should be judged based on their mental intentions to a notion that such mental intentions can be directly measured by potential future actions, with said future actions themselves being the arbiter of a person's judgment. Generally, the use of judgment is as a tool to prevent further ill deeds in others by whatever means are decided upon by the judgers. This is not to say it is usually an ideal practice, but we may nonetheless suppose that for oneself to judge a person might have a direct effect on that person's future actions, and such an effect would itself be only satisfactory in the first place if it meant preventing that person's destruction of Golden Scores, (or perhaps enabling their creation of Golden Scores). In such a case, we see that this model covers for both virtue ethics and consequentialist ethics, but also yields the wise insight that one ought not to judge a person they cannot control or change. 

\subsection{The Concept of Evil}

Let us now discuss how Golden ethics might parse ``good" versus ``evil". In order to assess such a thing, we are in need of exactly three Aethic degrees of freedom.
\begin{itemize}
    \item The person under analysis, (or alternatively a given abstract person given by a parent Aethus to a real person), who will be described by Aethus $A$, (explicitly relative to the Aethic reality given by Aethus $R$, which will therefore be incorporated into $A$ as needed). We will refer to this person as ``Person A" and their Aethus simply as $A$. 
    \item The particular Aethic attribute of Person A under analysis, being denoted by $\varphi$.
    \item A ``reference person", Person B, with Aethus $B$. Importantly, we necessitate---analogously to the first postulate of the Aethus, but running over ethics rather than reality---that no Person A has a particular well-defined goodness score: \textit{instead such may only be described relative to some Person B}.
    
    The closest we may get to such an ``objective goodness score" will be to take the extremal particular goodness score in a particular category against a collection of different possible Persons B, be it a demographic or Earth's entire population. 
\end{itemize}

% Anyways so the idea for this discussion here is that one's own choice of branding something as ethical or not in a given context is ITSELF a decision that could influence the Golden Point, and so should be assessed accordingly. The idea I want to emphasize though is that this is especially retrospectivaly-available structure to somebody who has already reached the Golden Point and is analyzing us after the fact---so this is how such an ``objective viewer" might compose an ethical framework with which to assess us as if we were their metaphorical first-person video game avatar of sorts or something! 
% One tangible example I want to provide about is the example of a malicious person inflicting harm on somebody else (for instance, shredding their college acceptance letter without telling them, etc, but maybe we come up with another example) which accidentally sets them on a different path which ultimately helps the Golden Point. So, the idea is that Automorality classifies the action as moral, but Relative Morality classifies it as malicious. Specifically, the malice class of it comes from supposing that in general, should we place this same person in that class of position of power again, and they would be expected to also do the destructive thing. So, even if it helped that one time, if we in general gave them that role, most of the time they would wreak havoc! 

We then define two parameters, being ``Automorality" versus ``Relative Morality". Such may be defined as follows.
\begin{itemize}
    \item \textbf{Automorality:}

    The automorality of Person A for attribute $\varphi$ quantifies the net increase in the natural logarithm of the Cosmic Nexus-probability of the Golden Point from Aethus $A$ to the complement of $A$ under $\varphi$, $A'_{\varphi}$.
    \begin{equation}
        \textnormal{AM} \! \left( A, \varphi \right) = \ln \! \left( \dfrac{\probP \left( G \mid A'_{\varphi}, N_0 \right)}{\probP \left( G \mid A, N_0 \right)} \right)
    \end{equation}
    This is so where $N_0$ is the Cosmic Nexus, $G$ is the Golden Aethus, and ``$\textnormal{AM} \! \left( A, \varphi \right)$" denotes automorality. 

    Intuitively, this denotes a person's strict behaviorist morality in terms of the direct aid they provide to the Golden Point probability. Remarkably, this will rank actions of wretched evolution as having positive automorality, because technically they aid the Golden Point, even if inadvertently. 
    
    \item \textbf{Relative Morality:}

    The morality of Person A for attribute $\varphi$ relative to Person B quantifies the net increase in the natural logarithm of the Cosmic Nexus-probability of the Golden Point from Aethus $B$ to the Aethus $B_{\phi, A}$ given by forming the Aethus which replaces $B$'s stated-attribute under $\phi$ with that of $A$.
    \begin{equation}
        \textnormal{RM}_{B} \! \left( A, \varphi \right) = \ln \! \left( \dfrac{\probP \left( G \mid B_{\phi, A}, N_0 \right)}{\probP \left( G \mid B, N_0 \right)} \right)
    \end{equation}
    This is so where $N_0$ is the Cosmic Nexus, $G$ is the Golden Aethus, and ``$\textnormal{RM}_{B} \! \left( A, \varphi \right)$" denotes relative morality. 

    Structurally, relative morality can also be defined as the automorality of Person B changing their own Aethus exclusively so that their stated-attribute under $\varphi$ reflects the stated-attribute which $A$ has under $\varphi$.

    Intuitively, this is meant to denote a person's more intentional morality, because it captures the ``inadvisable" aspect of immoral actions. That is, by testing the automorality of Person B directly if they were to become like Person A in said targeted way determined by $\varphi$. Accordingly, we have that events of wretched evolution would be properly categorized as relatively immoral to an agent who is meant to help the Golden Point by acting differently to that. 
\end{itemize}

We may now categorize the score of attribute $\varphi$ of Person A respecting to Person B with the following table. 

\begin{table}[htbp]
    \centering
    \textbf{\large {Good vs. Evil in Golden Ethics}} \\
    \vspace{0.5em}
    \begin{tabular}{|c|c|c|c|}
        \cline{2-4}
        \multicolumn{1}{c|}{} & \textbf{Autoimmoral} ($-1$) & \textbf{Autoneutral} ($0$) & \textbf{Automoral} ($+1$) \\
        \hline
        \textbf{Relatively Moral} ($+1$) & Biodepartric Evil & Ineffective & \textbf{Golden Good} \\
        \hline
        \textbf{Relatively Neutral} ($0$) & Departic Power & Neffective & Effective Power \\
        \hline
        \textbf{Relatively Immoral} ($-1$) & Departric Evil & Pissantry & Biotic Evil \\
        \hline
    \end{tabular}
    \caption{This is the table of how good versus evil may be generally regarded in Golden reasoning. Note that these concepts of good and evil ought not to taken as directly reflective of the standard societal perceptions, since of course that would result in an extreme loss of nuance. Instead, these terms have been given those names similarly to how Newton provided his quantity of force with the same name as the preexisting subjective quality of force.\\
    Note how the use of the world evil in some of these categories is meant to emphasize how the classical notion of evil cannot distinguish between these two states, even though in Golden reasoning we posit that such a distinction is highly important. }
    \label{tab:evil_types}
\end{table}

Interestingly, we may also now think to quantify the concept of ``power" in terms of this display as well.
\begin{equation}
    \textnormal{Power}_{B} \! \left( A, \varphi \right) = \left| {\dfrac{\textnormal{AM} \! \left( A, \varphi \right)}{\textnormal{RM}_{B} \! \left( A, \varphi \right)}} \right|
\end{equation}
That is, in the table, this is ``the absolute value of the $x$-coordinate divided by the absolute value of the $y$-coordinate, $\frac{\left| x \right|}{\left| y \right|}$". 

\subsection{Departric Invalidity}

Considering some of the core tenets to Golden reasoning, we might ask what they imply about a scenario in which the Departure Point occurs. Remembering how a Departure Point can be regarded as an event horizon between a biosphere and its chance of reaching a Golden Point, it then follows that the Golden Answer, in such a scenario, can no longer be applied to give meaningful answers about reality, because, blatantly, nothing than can occur has a correlation to the Golden Point anyway. As such, we can define a concept of ``Departric Invalidity" as follows.
\begin{definition}[Departric Invalidity]
    Once a biosphere has crossed its own Departure Point, all concepts of virtue or meaning within that biosphere cease to conceptually exist, similarly to what happens to time at the event horizon of a black hole, or to reality at the instance of an invalid Aethus. 
\end{definition}
There are many lenses through which we can analyze such a thing. Spiritually, we might suppose that it is somewhat similar to having one's Aethus be severed from the background flow of spirituality in the universe, like an end to our ``spiritual ventilation" from concepts like Brahman or the Tao. Or, on a more philosophical note, we might suppose that it removes all notions of duty or justice from reality, which are supplemented with anarchic equivalents. This includes concepts which we have a comfortable grasp on now, like the right to life. In a post-Departure scenario, we cannot even consider one's own right to life as being personally preferable over the alternative. Any lines we could draw which support this argument would have evaporated, so hence no basis could be drawn anyway for it. Stagnation of a pungent vanity would be left in place of intrinsic moral reasoning. 

Although this scenario is quite frightening, it also offers us a clear alternative derivation of Golden Ethics: that is, considering how the instance of the Departure Point deletes all morality anyway, it is therefore worth it at all costs to prevent this event. Let us present this in the form of the following thought experiment, which uses a foundational human instinct as a guide: self-preservation. Suppose there is a special amulet at the top of a skyscraper that you have to get a hold of in order to trigger the Golden Point. If you do not get it in the next hour, then the Departure Point will happen immediately afterwards. Also suppose you've looked at the stats, and have determined that you only have a $1\%$ chance of surviving all the way to the top of the building, as the danger level is extreme in this case. Also suppose that there is no one who is more skilled at this task who can fill in for you, so your options are to either take the climb, or accept that the Departure Point becomes an inevitability. You then ask yourself the following question: \textit{what is the best choice to take in terms of self-preservation}? Classically, we might think that the answer is obvious -- don't take the climb, because that gives a 100\% chance of survival rather than only a 1\% chance. However, when we consider Departric invalidity, we see a very different perspective. According to this principle, the state of the Departure Point occurring can be treated almost like an invalid Aethus, at least in terms of virtue. That is, whatever the chances of you physically surviving may be, you will lose any and all interpretation of such a statistic's value once the Departure Point occurs, so therefore it can be treated as indeterminate relative to another Departure Point Aethus, or alternatively as zero relative to your current Aethus. Therefore, you are valuing your life with $1\%$ effectiveness should you take the climb, but with $0\%$ effectiveness if you do not, so altogether we see that self-preservation\footnote{One might object to this by stating that self-preservation is actually a direct numerical result of probability, and hence is not actually invalidated in the event of the Departure Point. But, to counte-argue this, this is not quite so -- at least regarding the way we usually portray the concept of self-preservation, where it has many nuanced components thrown into the mix as well, (for example, it might extend to family, to beliefs, etcetera). Even a hint of nuance in a concept implies that that much of it is to blow a fuse of sorts in the event of the Departure Point, and hence will tend to cancel itself out regarding any semblance of continued ``existence." Hence, self-preservation is indeed no exception to Departric invalidity.}, (or perhaps self-preservation coupled with the continuation of a concept of self-preservation in the first place), is actually maximized by your climbing the building. 

Note how such a thing re-derives Golden Ethics -- that is, one ought to pick an option which is maximized in success probability given the Golden Point intersected with their current Aethus, rather than just their current Aethus alone. That is, the fundamental principle of Golden Ethics can be phrased through Departric invalidity as follows: ``\textit{pretend the Departure Point is an invalid Aethus, and then act in accordance with the highest probabilities}."

\subsection{Golden Choosing}

We can actually reformulate the fundamental principle of Golden Ethics using this application just now granted by Departric invalidity, that is, we decide that Aethus $X$ is more ``virtuous" than Aethus $Y$ as follows, where $A$ is one's own Aethus and $G$ is the Golden Aethus.
\begin{equation}
    \mathbf{M}_A \! \left( X \right) \geq \mathbf{M}_A \! \left( Y \right) \Leftrightarrow P \! \left( \frac{X}{Y} \middle| A + G \right) \geq P \! \left( \frac{Y}{X} \middle| A + G \right)
\end{equation}
This is so where $\frac{X}{Y}$ represents Aethic division, (being the removal of all attributes in $Y$ from $X$), and $X + Y$ represents Aethic intersection. Also, $\mathbf{M} \! \left( X \right)$ stands for the ``morality" of $X$, with the necessity of a subscript Aethus following from the first postulate of the Aethus. Put verbally, this says that \textit{one ought to pick a child Aethus as most moral which is likeliest given one's own Aethus together with the Golden Point}. Note how this is analogous to picking that attribute which makes the Golden Point likeliest, because it is again something of a compilation of the probability of the Golden Point given some task with the probability from your current Aethus of completing said task. 